% =====================================================================
% frontmatter/abstract.tex
% PT: Resumo em inglês (Abstract). Um único parágrafo, normalmente com
%     150 a 300 palavras (confirme o limite no regulamento em vigor).
%     Deve dizer: contexto, problema, abordagem, principais resultados e
%     conclusão. As palavras-chave (normalmente 3 a 6) vão no fim, na linha
%     "Keywords", separadas por ponto e vírgula.
% EN: Abstract in English. A single paragraph, usually 150 to 300 words
%     (check the limit in the current regulations). It should state:
%     context, problem, approach, main results and conclusion. The keywords
%     (usually 3 to 6) go at the end, on the "Keywords" line, separated by
%     semicolons.
% =====================================================================
\chapter*{Abstract}
\addcontentsline{toc}{chapter}{Abstract}

% PT: Texto de exemplo em inglês: substitua pelo seu resumo.
% EN: Example text in English: replace it with your abstract.
This is an example abstract. In a real dissertation, this paragraph
summarises the whole work for a reader who will not read anything else. It
starts with one or two sentences of context that explain why the topic
matters, then states the problem or gap the dissertation addresses. Next, it
describes the approach in concrete terms: the method, the data or case
study, and the tools that were used. It then reports the main results,
preferably with the key numbers, and closes with the main conclusion and
what it means for practice or for future research. The abstract should not
contain citations, figures, tables or undefined acronyms, and every claim in
it should be supported by the main text. Write it last, once the results and
conclusions are final, and keep it within the word limit set by the school.

\vspace{1em}
% PT: Palavras-chave em inglês (3 a 6), separadas por ponto e vírgula.
% EN: Keywords in English (3 to 6), separated by semicolons.
\noindent\textbf{Keywords:} Keyword One; Keyword Two; Keyword Three; Keyword Four
